\section{结论}

本文针对全景语义分割中的训练效率和边界不连续问题，提出了一种高效的记忆优化框架。通过记忆Token压缩(MTC)、稀疏记忆检索(SMR)和无缝循环记忆(SCM)三项创新，我们实现了精度与效率的良好平衡。在Cityscapes到DensePASS的无监督域适应任务中，MTC+SMR组合实现8倍Token压缩和55.8\%注意力时间降低，仅损失0.24\%的mIoU。SCM在推理阶段提升整体精度0.33\%，交通标志类别提升5.04\%，验证了边界连续性建模的有效性。5个随机种子的多次实验表明我们的方法具有良好的稳定性和可重复性。

未来工作包括：(1)在更大规模数据集(如完整的Cityscapes 2975张图像)上验证方法的扩展性；(2)探索MTC和SCM与其他域适应技术(如对抗训练、原型对齐)的协同效果；(3)在嵌入式设备(如Jetson系列)上部署并测试实际推理性能；(4)扩展到多任务场景(如全景分割与深度估计的联合学习)。

\section*{致谢}

本研究得到国家自然科学基金(项目编号: XXXXXXXX)的支持。感谢AutoDL平台提供的GPU计算资源。
